\section{Mechanism-Ablation and Rescue Analysis}
\label{app:mechanism-analysis}

This appendix defines the mechanism variants used in
Section~\ref{sec:eval:mechanism} and the conservative trace-level criterion
behind the reported pre-handoff rescue statistic.
All mechanism results and the rescue decomposition use the same 500-request
Burst $12\times$ execution; the decomposition is not a separate experimental
campaign.


% --------------------------------------------------------------------------
\subsection{Mechanism-Ablation Definitions}
\label{sec:appendix:mechanism-ablations}

The ablations remove one control mechanism at a time while retaining the
remaining RAVEL lifecycle wherever possible.
Table~\ref{tab:appendix:mechanism-ablations} summarizes the comparison.

\begin{table}[t]
    \centering
    \small
    \caption{Mechanism variants used in the Burst $12\times$ ablation.}
    \label{tab:appendix:mechanism-ablations}
    \setlength{\tabcolsep}{3.5pt}
    \renewcommand{\arraystretch}{1.10}
    \begin{tabular}{@{}lccc@{}}
        \toprule
        \textbf{Variant} &
        \makecell{\textbf{Prospective}\\\textbf{accounting}} &
        \makecell{\textbf{Pre-handoff}\\\textbf{reassignment}} &
        \makecell{\textbf{Local TBT}\\\textbf{protection}} \\
        \midrule
        RAVEL          & Yes & Yes & Yes \\
        NoAccounting   & No  & Yes & Yes \\
        FixedPlacement & Yes & No  & Yes \\
        NoTBTGuard     & Yes & Yes & No  \\
        \bottomrule
    \end{tabular}
\end{table}

\paragraph{NoAccounting.}
This variant retains replica-specific SLO quotes, the pre-handoff replanning
path, the ownership/commitment boundary, and engine-local TBT protection, but
suppresses the normal prospective representation of Router-owned tentative
assignments in the capacity state consumed by subsequent placement decisions.
Live engine pending and running state remains visible.
The variant therefore isolates the value of making uncommitted demand visible
before handoff.

\paragraph{FixedPlacement.}
This variant performs RAVEL's initial SLO-aware placement and retains
prospective accounting, dispatch/handoff semantics, and engine-local TBT
protection, but disables ordinary pre-handoff reassignment after the initial
choice.
It isolates the value of keeping Router-owned placement revisable while holding
tentative-demand visibility fixed.

\paragraph{NoTBTGuard.}
This variant retains global quotes, prospective accounting, placement, and
pre-handoff replanning, but disables the engine-local Prefill control and
priority behavior used to protect active LATENCY Decode requests.
It isolates the local mechanism that preserves the TBT assumptions used by
replica-level feasibility decisions.


% --------------------------------------------------------------------------
\subsection{Confirmed Pre-Handoff Rescue Criterion}
\label{sec:appendix:rescue-criterion}

The trace diagnostic asks whether a miss under a frozen pre-handoff placement
can be conservatively attributed to an allocation that the same planning
episode later repairs.
We call the affected request the \emph{victim} and an earlier-arrived request
that vacates the region obtained by the victim in the same plan the
\emph{blocker}.

A frozen-placement miss is classified as a \emph{confirmed pre-handoff rescue}
only when all of the following conditions hold:
\begin{enumerate}
    \item Under the incumbent allocation frozen at the same pre-handoff planning
          episode, the victim is predicted infeasible.
    \item After the cohort reassignment produced by that episode, the victim is
          predicted feasible.
    \item The victim subsequently satisfies its realized request-level SLO.
    \item At least one blocker arrived earlier, occupied the region that the
          victim obtains, and is moved away from that region by the same plan.
\end{enumerate}
The criterion intentionally requires both a decision-level change in predicted
feasibility and a successful realized outcome; cases that cannot establish all
four conditions remain unresolved.

\paragraph{Feasible-set reconstruction.}
The diagnostic must distinguish intrinsic regional infeasibility from
congestion created by the very Router cohort being analyzed.
For each relevant pre-handoff snapshot, we therefore retain materialized engine
load but remove the current Router cohort from prospective placement state.
We then test the affected request individually against each candidate region
using the same causal quote machinery.
This reconstruction yields the diagnostic feasible-region set without treating
blocker-induced prospective congestion as an intrinsic property of the region.
No future arrivals or realized output lengths are introduced into the quote.

\paragraph{Interpretation boundary.}
Under this conservative criterion, at least 41.25\% of misses under frozen
pre-handoff placement are confirmed redirect-rescuable.
The remaining at most 58.75\% are \emph{unresolved} by this diagnostic; they
should not be interpreted as proven capacity shortages or irrecoverable
failures.
The statistic decomposes the same Burst $12\times$ mechanism run used for the
ablation rather than supplying an independent experimental sample.
It provides trace-level evidence for the mechanism behind the
FixedPlacement gap, not a counterfactual guarantee that every classified miss
would be repaired under arbitrary execution.

\paragraph{Trace collection.}
The diagnostic records the planning episode, incumbent and revised ownership,
predicted feasibility before and after reassignment, and subsequent realized
SLO outcome.
These records are buffered during execution and persisted at the end of the
run so that instrumentation does not introduce synchronous logging work into
the scheduling path.


% ==========================================================================
% Appendix C
% ==========================================================================
